\documentclass{article}

% -----------------------------
% Packages
% -----------------------------
\usepackage{fancyhdr}
\usepackage{extramarks}
\usepackage{amsmath}
\usepackage{amsthm}
\usepackage{amsfonts}
\usepackage{tikz}
\usepackage[plain]{algorithm}
\usepackage{algpseudocode}
\usepackage{amssymb}
\usepackage{enumitem}

\usetikzlibrary{automata,positioning}

% -----------------------------
% Basic Document Settings
% -----------------------------
\topmargin=-0.45in
\evensidemargin=0in
\oddsidemargin=0in
\textwidth=6.5in
\textheight=9.0in
\headsep=0.25in

\linespread{1.1}

\pagestyle{fancy}
\lhead{\hmwkAuthorName}
\chead{\hmwkClass: \hmwkTitle}
\rhead{}
\lfoot{\lastxmark}
\cfoot{\thepage}

\renewcommand\headrulewidth{0.4pt}
\renewcommand\footrulewidth{0.4pt}

\setlength\parindent{0pt}

% -----------------------------
% Problem Section Helpers
% -----------------------------
\newcommand{\enterProblemHeader}[1]{
  \nobreak\extramarks{}{Problem #1 continued on next page\ldots}\nobreak{}
  \nobreak\extramarks{Problem #1 (continued)}{Problem #1 continued on next page\ldots}\nobreak{}
}

\newcommand{\exitProblemHeader}[1]{
  \nobreak\extramarks{Problem #1 (continued)}{Problem #1 continued on next page\ldots}\nobreak{}
  \nobreak\extramarks{Problem #1}{}\nobreak{}
}

\setcounter{secnumdepth}{0}
\newcounter{partCounter}
\newcounter{homeworkProblemCounter}
\setcounter{homeworkProblemCounter}{1}
\nobreak\extramarks{Problem \arabic{homeworkProblemCounter}}{}\nobreak{}

% -----------------------------
% Theorem env (optional)
% -----------------------------
\newtheorem*{theorem}{Theorem}

% -----------------------------
% Homework Problem Environment
%   Usage:
%   \begin{homeworkProblem}            % auto-number
%     % Your answer here
%   \end{homeworkProblem}
%
%   \begin{homeworkProblem}[1.2.3]     % label problem "1.2.3" (e.g. textbook numbering)
%     % Your answer here
%   \end{homeworkProblem}
% -----------------------------
\newenvironment{homeworkProblem}[1][]{
  \def\hwProblemLabel{#1}
  \ifx\hwProblemLabel\empty
    \edef\hwProblemLabel{\arabic{homeworkProblemCounter}}
    \stepcounter{homeworkProblemCounter}
  \fi
  \section{Problem \hwProblemLabel}
  \setcounter{partCounter}{1}
  \enterProblemHeader{\hwProblemLabel}
}{
  \exitProblemHeader{\hwProblemLabel}
}

% -----------------------------
% Homework Details (EDIT THESE)
% -----------------------------
\newcommand{\hmwkTitle}{Homework \#5}
\newcommand{\hmwkClass}{Math 3320}
\newcommand{\hmwkAuthorName}{\textbf{Joseph Quinn}}

% -----------------------------
% Title
% -----------------------------
\title{
  \vspace{2in}
  \textmd{\textbf{\hmwkClass:\ \hmwkTitle}}\\
  \vspace{3in}
}
\author{\hmwkAuthorName}
\date{}

% -----------------------------
% Convenience Macros (optional)
% -----------------------------
\renewcommand{\part}[1]{\textbf{\large Part \Alph{partCounter}}\stepcounter{partCounter}\\}

% Algorithms
\newcommand{\alg}[1]{\textsc{\bfseries \footnotesize #1}}

% Calculus
\newcommand{\deriv}[1]{\frac{\mathrm{d}}{\mathrm{d}x} (#1)}
\newcommand{\pderiv}[2]{\frac{\partial}{\partial #1} (#2)}
\newcommand{\dx}{\mathrm{d}x}

% Statistics
\newcommand{\E}{\mathrm{E}}
\newcommand{\Var}{\mathrm{Var}}
\newcommand{\Cov}{\mathrm{Cov}}
\newcommand{\Bias}{\mathrm{Bias}}

% Proof step (for aligned derivations)
\newcommand{\step}[2]{& #1 & & \text{#2} \\}

% Solution header (optional)
\newcommand{\solution}{\textbf{\large Solution}}

% -----------------------------
% Document
% -----------------------------
\begin{document}

\maketitle
\pagebreak


\begin{homeworkProblem}[2.6.7]

	\begin{enumerate}
		\item[(b)]
		      \[
			      G =
			      \begin{pmatrix}
				      1 & 1 & 1 & 1 & 1 & 1 & 0 & 0 \\
				      0 & 0 & 1 & 1 & 0 & 0 & 1 & 1 \\
				      0 & 0 & 0 & 0 & 1 & 1 & 1 & 1 \\
				      0 & 0 & 0 & 0 & 0 & 0 & 1 & 1
			      \end{pmatrix}
		      \]

		      Since the four rows are linearly independent,
		      \(k=4\). And
		      \(00000011\) is the min weight codeword being \(2\).

		      \[
			      (n,k,d)=(8,4,2)
		      \]

		\item[(d)]
		      \[
			      G =
			      \begin{pmatrix}
				      1 & 0 & 1 & 0 & 1 \\
				      0 & 1 & 0 & 1 & 0 \\
				      0 & 0 & 0 & 1 & 1
			      \end{pmatrix}
		      \]

		      The three rows are linearly independent,
		      so \(k=3\). The minimum nonzero codeword
		      weight is \(2\) from \(01010\).

		      \[
			      (n,k,d)=(5,3,2)
		      \]

		\item[(e)]
		      \[
			      G =
			      \begin{pmatrix}
				      1 & 0 & 1 & 0 \\
				      0 & 1 & 0 & 1
			      \end{pmatrix}
		      \]

		      Since \(1111=1010+0101\), the third
		      vector is dependent on the first two.
		      Thus \(k=2\).

		      And the smallest weight from the matrix is $2$

		      \[
			      (n,k,d)=(4,2,2)
		      \]
	\end{enumerate}


\end{homeworkProblem}
\begin{homeworkProblem}[2.6.12]

	\[
		G=
		\begin{pmatrix}
			1 & 0 & 0 & 0 & 1 & 1 & 1 \\
			0 & 1 & 0 & 0 & 1 & 1 & 0 \\
			0 & 0 & 1 & 0 & 1 & 0 & 1 \\
			0 & 0 & 0 & 1 & 0 & 1 & 1
		\end{pmatrix}
	\]

	Let \(g_1,g_2,g_3,g_4\) denote the rows of \(G\).
	For each message \(m\), the codeword is \(c=mG\).
	\begin{enumerate}
		\item[(A)] $0000G = 0g_1+0g_2+0g_3+0g_4 = 0000000$
		\item[(B)] $1000G = g_1 = 1000111$
		\item[(C)] $0100G = g_2 = 0100110$
		\item[(D)] $0010G = g_3 = 0010101$
		\item[(E)] $0001G = g_4 = 0001011$
		\item[(F)] $1100G = g_1+g_2 = 1000111+0100110 = 1100001$
		\item[(G)] $1010G = g_1+g_3 = 1000111+0010101 = 1010010$
		\item[(H)] $1001G = g_1+g_4 = 1000111+0001011 = 1001100$
		\item[(I)] $0110G = g_2+g_3 = 0100110+0010101 = 0110011$
		\item[(J)] $0101G = g_2+g_4 = 0100110+0001011 = 0101101$
		\item[(K)] $0011G = g_3+g_4 = 0010101+0001011 = 0011110$
		\item[(L)] $1110G = g_1+g_2+g_3 = 1000111+0100110+0010101 = 1110100$
		\item[(M)] $1101G = g_1+g_2+g_4 = 1000111+0100110+0001011 = 1101010$
		\item[(N)] $1011G = g_1+g_3+g_4 = 1000111+0010101+0001011 = 1011001$
		\item[(O)] $0111G = g_2+g_3+g_4 = 0100110+0010101+0001011 = 0111000$
		\item[(P)] $1111G = g_1+g_2+g_3+g_4 = 1000111+0100110+0010101+0001011 = 1111111$
	\end{enumerate}


	we can encode the message Call Home Bama as:
	\[
		\begin{aligned}
			\text{CALL} & \rightarrow
			0100110\ 0000000\ 1110100\ 1110100 \\
			\text{HOME} & \rightarrow
			1001100\ 0111000\ 1101010\ 0001011 \\
			\text{BAMA} & \rightarrow
			1000111\ 0000000\ 1101010\ 0000000
		\end{aligned}
	\]
\end{homeworkProblem}
\begin{homeworkProblem}[2.7.5]


	\begin{enumerate}
		\item[(c)]
		      \[
			      C=\{0000000000,1111100000,0000011111,1111111111\}.
		      \]
		      the zero vector makes this dependent as well as,
		      \[
			      1111111111=1111100000+0000011111.
		      \]

		      Meaning the largest set of independent words is the following,
		      \[
			      G=
			      \begin{pmatrix}
				      1 & 1 & 1 & 1 & 1 & 0 & 0 & 0 & 0 & 0 \\
				      0 & 0 & 0 & 0 & 0 & 1 & 1 & 1 & 1 & 1
			      \end{pmatrix}.
		      \]

		      we have two linearly independent rows so $k=2$ and $n-k=8$.

		      To put $G$ in standard form, move
		      column 6 to position 2.

		      \[
			      G'=
			      \left(
			      \begin{array}{cc|cccccccc}
				      1 & 0 & 1 & 1 & 1 & 1 & 0 & 0 & 0 & 0 \\
				      0 & 1 & 0 & 0 & 0 & 0 & 1 & 1 & 1 & 1
			      \end{array}
			      \right)
		      \]

		      \[
			      X=
			      \begin{pmatrix}
				      1 & 1 & 1 & 1 & 0 & 0 & 0 & 0 \\
				      0 & 0 & 0 & 0 & 1 & 1 & 1 & 1
			      \end{pmatrix}.
		      \]

		      We construct
		      \[
			      H'=
			      \begin{pmatrix}X\\I_8\end{pmatrix}
			      =
			      \begin{pmatrix}
				      1 & 1 & 1 & 1 & 0 & 0 & 0 & 0 \\
				      0 & 0 & 0 & 0 & 1 & 1 & 1 & 1 \\
				      1 & 0 & 0 & 0 & 0 & 0 & 0 & 0 \\
				      0 & 1 & 0 & 0 & 0 & 0 & 0 & 0 \\
				      0 & 0 & 1 & 0 & 0 & 0 & 0 & 0 \\
				      0 & 0 & 0 & 1 & 0 & 0 & 0 & 0 \\
				      0 & 0 & 0 & 0 & 1 & 0 & 0 & 0 \\
				      0 & 0 & 0 & 0 & 0 & 1 & 0 & 0 \\
				      0 & 0 & 0 & 0 & 0 & 0 & 1 & 0 \\
				      0 & 0 & 0 & 0 & 0 & 0 & 0 & 1
			      \end{pmatrix}.
		      \]

		      We must restore the original order now resulting in


		      \[
			      H=
			      \begin{pmatrix}
				      1 & 1 & 1 & 1 & 0 & 0 & 0 & 0 \\
				      1 & 0 & 0 & 0 & 0 & 0 & 0 & 0 \\
				      0 & 1 & 0 & 0 & 0 & 0 & 0 & 0 \\
				      0 & 0 & 1 & 0 & 0 & 0 & 0 & 0 \\
				      0 & 0 & 0 & 1 & 0 & 0 & 0 & 0 \\
				      0 & 0 & 0 & 0 & 1 & 1 & 1 & 1 \\
				      0 & 0 & 0 & 0 & 1 & 0 & 0 & 0 \\
				      0 & 0 & 0 & 0 & 0 & 1 & 0 & 0 \\
				      0 & 0 & 0 & 0 & 0 & 0 & 1 & 0 \\
				      0 & 0 & 0 & 0 & 0 & 0 & 0 & 1
			      \end{pmatrix}.
		      \]

		\item[(d)]
		      \[
			      C=\{11111111,11110000,11001100,10101010\}.
		      \]

		      To setup the Generator matrix lets write this in matrix form
		      \[
			      G=
			      \begin{pmatrix}
				      1 & 1 & 1 & 1 & 1 & 1 & 1 & 1 \\
				      1 & 1 & 1 & 1 & 0 & 0 & 0 & 0 \\
				      1 & 1 & 0 & 0 & 1 & 1 & 0 & 0 \\
				      1 & 0 & 1 & 0 & 1 & 0 & 1 & 0
			      \end{pmatrix}.
		      \]

		      And now row reduce to RREF

		      $R_2\leftarrow R_2+R_1$,
		      $R_3\leftarrow R_3+R_1$,
		      and $R_4\leftarrow R_4+R_1$.

		      \[
			      G\sim
			      \begin{pmatrix}
				      1 & 1 & 1 & 1 & 1 & 1 & 1 & 1 \\
				      0 & 0 & 0 & 0 & 1 & 1 & 1 & 1 \\
				      0 & 0 & 1 & 1 & 0 & 0 & 1 & 1 \\
				      0 & 1 & 0 & 1 & 0 & 1 & 0 & 1
			      \end{pmatrix}.
		      \]

		      Swap $R_2$ and $R_4$:

		      \[
			      \sim
			      \begin{pmatrix}
				      1 & 1 & 1 & 1 & 1 & 1 & 1 & 1 \\
				      0 & 1 & 0 & 1 & 0 & 1 & 0 & 1 \\
				      0 & 0 & 1 & 1 & 0 & 0 & 1 & 1 \\
				      0 & 0 & 0 & 0 & 1 & 1 & 1 & 1
			      \end{pmatrix}.
		      \]

		      $R_1\leftarrow R_1+R_2+R_3+R_4$.

		      \[
			      \sim
			      \begin{pmatrix}
				      1 & 0 & 0 & 1 & 0 & 1 & 1 & 0 \\
				      0 & 1 & 0 & 1 & 0 & 1 & 0 & 1 \\
				      0 & 0 & 1 & 1 & 0 & 0 & 1 & 1 \\
				      0 & 0 & 0 & 0 & 1 & 1 & 1 & 1
			      \end{pmatrix}.
		      \]

		      There are four pivots, so the
		      generators are linearly independent.
		      Thus $k=4$ and $n-k=4$.

		      Swap columns 4 and 5 to put in standard form:

		      \[
			      G'=
			      \left(
			      \begin{array}{cccc|cccc}
				      1 & 0 & 0 & 0 & 1 & 1 & 1 & 0 \\
				      0 & 1 & 0 & 0 & 1 & 1 & 0 & 1 \\
				      0 & 0 & 1 & 0 & 1 & 0 & 1 & 1 \\
				      0 & 0 & 0 & 1 & 0 & 1 & 1 & 1
			      \end{array}
			      \right).
		      \]
		      \[
			      X=
			      \begin{pmatrix}
				      1 & 1 & 1 & 0 \\
				      1 & 1 & 0 & 1 \\
				      1 & 0 & 1 & 1 \\
				      0 & 1 & 1 & 1
			      \end{pmatrix}.
		      \]

		      and begin to formulate our parity check matrix as,
		      \[
			      H'=
			      \begin{pmatrix}X\\I_4\end{pmatrix}
			      =
			      \begin{pmatrix}
				      1 & 1 & 1 & 0 \\
				      1 & 1 & 0 & 1 \\
				      1 & 0 & 1 & 1 \\
				      0 & 1 & 1 & 1 \\
				      1 & 0 & 0 & 0 \\
				      0 & 1 & 0 & 0 \\
				      0 & 0 & 1 & 0 \\
				      0 & 0 & 0 & 1
			      \end{pmatrix}.
		      \]

		      Now we have to restore $H'$ to its original order.

		      \[
			      H=
			      \begin{pmatrix}
				      1 & 1 & 1 & 0 \\
				      1 & 1 & 0 & 1 \\
				      1 & 0 & 1 & 1 \\
				      1 & 0 & 0 & 0 \\
				      0 & 1 & 1 & 1 \\
				      0 & 1 & 0 & 0 \\
				      0 & 0 & 1 & 0 \\
				      0 & 0 & 0 & 1
			      \end{pmatrix}.
		      \]

		\item[(f)]
		      \[
			      C=\{100100100,010010010,001001001,111111111\}.
		      \]

		      Now lets create the matrix representation as,

		      \[
			      G=
			      \begin{pmatrix}
				      1 & 0 & 0 & 1 & 0 & 0 & 1 & 0 & 0 \\
				      0 & 1 & 0 & 0 & 1 & 0 & 0 & 1 & 0 \\
				      0 & 0 & 1 & 0 & 0 & 1 & 0 & 0 & 1 \\
				      1 & 1 & 1 & 1 & 1 & 1 & 1 & 1 & 1
			      \end{pmatrix}.
		      \]

		      \[
			      R_4=R_1+R_2+R_3.
		      \]

		      $R_4\leftarrow R_4+R_1+R_2+R_3$:

		      \[
			      G\sim
			      \begin{pmatrix}
				      1 & 0 & 0 & 1 & 0 & 0 & 1 & 0 & 0 \\
				      0 & 1 & 0 & 0 & 1 & 0 & 0 & 1 & 0 \\
				      0 & 0 & 1 & 0 & 0 & 1 & 0 & 0 & 1 \\
				      0 & 0 & 0 & 0 & 0 & 0 & 0 & 0 & 0
			      \end{pmatrix}.
		      \]

		      drop the zero row:

		      \[
			      G=
			      \left(
			      \begin{array}{ccc|cccccc}
				      1 & 0 & 0 & 1 & 0 & 0 & 1 & 0 & 0 \\
				      0 & 1 & 0 & 0 & 1 & 0 & 0 & 1 & 0 \\
				      0 & 0 & 1 & 0 & 0 & 1 & 0 & 0 & 1
			      \end{array}
			      \right).
		      \]

		      This is already in standard form and we have $k=3$ and $n-k=6$.

		      \[
			      X=
			      \begin{pmatrix}
				      1 & 0 & 0 & 1 & 0 & 0 \\
				      0 & 1 & 0 & 0 & 1 & 0 \\
				      0 & 0 & 1 & 0 & 0 & 1
			      \end{pmatrix}.
		      \]

		      \[
			      H=
			      \begin{pmatrix}X\\I_6\end{pmatrix}
			      =
			      \begin{pmatrix}
				      1 & 0 & 0 & 1 & 0 & 0 \\
				      0 & 1 & 0 & 0 & 1 & 0 \\
				      0 & 0 & 1 & 0 & 0 & 1 \\
				      1 & 0 & 0 & 0 & 0 & 0 \\
				      0 & 1 & 0 & 0 & 0 & 0 \\
				      0 & 0 & 1 & 0 & 0 & 0 \\
				      0 & 0 & 0 & 1 & 0 & 0 \\
				      0 & 0 & 0 & 0 & 1 & 0 \\
				      0 & 0 & 0 & 0 & 0 & 1
			      \end{pmatrix}.
		      \]

	\end{enumerate}
\end{homeworkProblem}

\begin{homeworkProblem}[2.7.9]

	\begin{enumerate}
		\item[(c)]
		      \[
			      H=
			      \begin{pmatrix}
				      1 & 1 & 1 \\
				      1 & 1 & 0 \\
				      1 & 0 & 1 \\
				      0 & 1 & 1 \\
				      1 & 0 & 0 \\
				      0 & 1 & 0 \\
				      0 & 0 & 1
			      \end{pmatrix}.
		      \]

		      \textbf{(i)} A generator matrix for $C^\perp$ is simply $H^T$:

		      \[
			      G_{C^\perp}=H^T=
			      \begin{pmatrix}
				      1 & 1 & 1 & 0 & 1 & 0 & 0 \\
				      1 & 1 & 0 & 1 & 0 & 1 & 0 \\
				      1 & 0 & 1 & 1 & 0 & 0 & 1
			      \end{pmatrix}.
		      \]

		      \textbf{(ii)} To get a generator matrix for $C$, we need to perform the algo

		      \[
			      G'_{C^\perp}=
			      \left(
			      \begin{array}{ccc|cccc}
				      1 & 0 & 0 & 1 & 1 & 1 & 0 \\
				      0 & 1 & 0 & 1 & 1 & 0 & 1 \\
				      0 & 0 & 1 & 1 & 0 & 1 & 1
			      \end{array}
			      \right)
		      \]

		      \[
			      X=
			      \begin{pmatrix}
				      1 & 1 & 1 & 0 \\
				      1 & 1 & 0 & 1 \\
				      1 & 0 & 1 & 1
			      \end{pmatrix}.
		      \]

		      So we can begin to form $H'_{C^\perp}=\begin{pmatrix}X\\I_4\end{pmatrix}$,

		      and we need undoing the column permutation restores the original order.

		      \[
			      H_{C^\perp}=
			      \begin{pmatrix}
				      1 & 0 & 0 & 0 \\
				      0 & 1 & 0 & 0 \\
				      0 & 0 & 1 & 0 \\
				      0 & 0 & 0 & 1 \\
				      1 & 1 & 1 & 0 \\
				      1 & 1 & 0 & 1 \\
				      1 & 0 & 1 & 1
			      \end{pmatrix}.
		      \]

		      Finally, transpose to obtain a generator matrix for $C$:

		      \[
			      G_C=H_{C^\perp}^T=
			      \begin{pmatrix}
				      1 & 0 & 0 & 0 & 1 & 1 & 1 \\
				      0 & 1 & 0 & 0 & 1 & 1 & 0 \\
				      0 & 0 & 1 & 0 & 1 & 0 & 1 \\
				      0 & 0 & 0 & 1 & 0 & 1 & 1
			      \end{pmatrix}.
		      \]
	\end{enumerate}

\end{homeworkProblem}
\begin{homeworkProblem}[2.8.1]


	\[
		G=\begin{pmatrix}
			1 & 0 & 0 \\
			0 & 0 & 1
		\end{pmatrix}.
	\]

	The code generated by $G$ is

	\[
		C=\{000,100,001,101\}.
	\]

	Since $k=2$, any two linearly independent codewords form a generator matrix.

	The nonzero codewords are

	\[
		100,\quad 001,\quad 101.
	\]

	And order matters so we can write the other five generator matrices as

	\[
		G_1=\begin{pmatrix}
			0 & 0 & 1 \\
			1 & 0 & 0
		\end{pmatrix},
		\qquad
		G_2=\begin{pmatrix}
			1 & 0 & 0 \\
			1 & 0 & 1
		\end{pmatrix}
	\]

	\[
		G_3=\begin{pmatrix}
			1 & 0 & 1 \\
			1 & 0 & 0
		\end{pmatrix},
		\qquad
		G_4=\begin{pmatrix}
			0 & 0 & 1 \\
			1 & 0 & 1
		\end{pmatrix}
	\]

	\[
		G_5=\begin{pmatrix}
			1 & 0 & 1 \\
			0 & 0 & 1
		\end{pmatrix}.
	\]

\end{homeworkProblem}
\begin{homeworkProblem}[2.8.12]

	\[
		H_C=
		\begin{pmatrix}
			1 & 0 & 0 \\
			1 & 1 & 1 \\
			0 & 1 & 0 \\
			1 & 1 & 0 \\
			1 & 0 & 1 \\
			0 & 0 & 1 \\
			0 & 1 & 1
		\end{pmatrix}.
	\]

	We want to transpose $H_C$ to obtain a
	generator matrix for $C^\perp$:

	\[
		G_{C^\perp}=H_C^T=
		\begin{pmatrix}
			1 & 1 & 0 & 1 & 1 & 0 & 0 \\
			0 & 1 & 1 & 1 & 0 & 0 & 1 \\
			0 & 1 & 0 & 0 & 1 & 1 & 1
		\end{pmatrix}.
	\]


	Now we rearrange the columns into the order
	$(1,3,6,2,4,5,7)$ to form identity matrix

	\[
		G'_{C^\perp}=
		\left(
		\begin{array}{ccc|cccc}
			1 & 0 & 0 & 1 & 1 & 1 & 0 \\
			0 & 1 & 0 & 1 & 1 & 0 & 1 \\
			0 & 0 & 1 & 1 & 0 & 1 & 1
		\end{array}
		\right)
	\]

	\[
		X=
		\begin{pmatrix}
			1 & 1 & 1 & 0 \\
			1 & 1 & 0 & 1 \\
			1 & 0 & 1 & 1
		\end{pmatrix}.
	\]

	Now we can begin to write the parity check matrix.

	\[
		H'_{C^\perp}=
		\begin{pmatrix}
			1 & 1 & 1 & 0 \\
			1 & 1 & 0 & 1 \\
			1 & 0 & 1 & 1 \\
			1 & 0 & 0 & 0 \\
			0 & 1 & 0 & 0 \\
			0 & 0 & 1 & 0 \\
			0 & 0 & 0 & 1
		\end{pmatrix}.
	\]

	Because we want our final code to be in standard form and equivilent to $C$
	we want to moe hte last 4 rows to the front so after transpose its in standard form.

	\[
		H''_{C^\perp}=
		\left(
		\begin{array}{cccc}
			1 & 0 & 0 & 0 \\
			0 & 1 & 0 & 0 \\
			0 & 0 & 1 & 0 \\
			0 & 0 & 0 & 1 \\ \hline
			1 & 1 & 1 & 0 \\
			1 & 1 & 0 & 1 \\
			1 & 0 & 1 & 1
		\end{array}
		\right)
		=
		\begin{pmatrix}
			I_4 \\
			X
		\end{pmatrix}.
	\]

	Finally, transpose to obtain the
	generator matrix in standard form:

	\[
		G'=(H''_{C^\perp})^T=
		\left(
		\begin{array}{cccc|ccc}
			1 & 0 & 0 & 0 & 1 & 1 & 1 \\
			0 & 1 & 0 & 0 & 1 & 1 & 0 \\
			0 & 0 & 1 & 0 & 1 & 0 & 1 \\
			0 & 0 & 0 & 1 & 0 & 1 & 1
		\end{array}
		\right).
	\]

	The resulting code $C'$ is equivalent to $C$ and in standard form
\end{homeworkProblem}
\begin{homeworkProblem}[2.9.3]

	\[
		H_C=
		\begin{pmatrix}
			1 & 1 & 0 \\
			0 & 1 & 1 \\
			1 & 0 & 0 \\
			0 & 1 & 0 \\
			0 & 0 & 1
		\end{pmatrix}.
	\]

	First, transpose $H_C$ to obtain a generator
	matrix for $C^\perp$:

	\[
		G_{C^\perp}=H_C^T=
		\begin{pmatrix}
			1 & 0 & 1 & 0 & 0 \\
			1 & 1 & 0 & 1 & 0 \\
			0 & 1 & 0 & 0 & 1
		\end{pmatrix}.
	\]

	Notice that columns $3$, $4$, and $5$
	form the identity matrix $I_3$.

	Rearrange the columns into the order
	$(3,4,5,1,2)$.

	\[
		G'_{C^\perp}=
		\left(
		\begin{array}{ccc|cc}
			1 & 0 & 0 & 1 & 0 \\
			0 & 1 & 0 & 1 & 1 \\
			0 & 0 & 1 & 0 & 1
		\end{array}
		\right)
	\]
	\[
		X=
		\begin{pmatrix}
			1 & 0 \\
			1 & 1 \\
			0 & 1
		\end{pmatrix}.
	\]
	Now we can write the parity matrix as,

	\[
		H'_{C^\perp}=
		\begin{pmatrix}
			X \\
			I_2
		\end{pmatrix}.
	\]

	Substituting gives

	\[
		H'_{C^\perp}=
		\begin{pmatrix}
			1 & 0 \\
			1 & 1 \\
			0 & 1 \\
			1 & 0 \\
			0 & 1
		\end{pmatrix}.
	\]

	And we must undo our column permulation $(3,4,5,1,2)$ giving
	\[
		H_{C^\perp}=
		\begin{pmatrix}
			1 & 0 \\
			0 & 1 \\
			1 & 0 \\
			1 & 1 \\
			0 & 1
		\end{pmatrix}.
	\]

	Now if we transpose we obtain a generator
	matrix for $C$:

	\[
		G_C=H_{C^\perp}^T=
		\begin{pmatrix}
			1 & 0 & 1 & 1 & 0 \\
			0 & 1 & 0 & 1 & 1
		\end{pmatrix}.
	\]

	Now generate every codeword using $c=uG_C$.

	\[
		\begin{aligned}
			00G_C & = 0g_1+0g_2 = 00000 \\
			10G_C & = g_1 = 10110       \\
			01G_C & = g_2 = 01011       \\
			11G_C & = g_1+g_2 = 11101
		\end{aligned}
	\]
	Meaning the code generated by this matrix is,
	\[
		C=\{00000,10110,01011,11101\}.
	\]

	And with this we see that the minimum weight of a non zero codeword is the choice of any of the non zero codewords, being $wt(10110)=3$
\end{homeworkProblem}
\begin{homeworkProblem}[2.10.8]
	\[
		H_C=
		\begin{pmatrix}
			1 & 1 & 1 \\
			1 & 1 & 0 \\
			1 & 0 & 1 \\
			0 & 1 & 1 \\
			1 & 0 & 0 \\
			0 & 1 & 0 \\
			0 & 0 & 1
		\end{pmatrix}.
	\]

	First, transpose $H_C$ to obtain a generator
	matrix for $C^\perp$:

	\[
		G_{C^\perp}=H_C^T=
		\begin{pmatrix}
			1 & 1 & 1 & 0 & 1 & 0 & 0 \\
			1 & 1 & 0 & 1 & 0 & 1 & 0 \\
			1 & 0 & 1 & 1 & 0 & 0 & 1
		\end{pmatrix}.
	\]

	Now we can rearrange the columns into the order
	$(5,6,7,1,2,3,4)$ to form identity matrix

	\[
		G'_{C^\perp}=
		\left(
		\begin{array}{ccc|cccc}
			1 & 0 & 0 & 1 & 1 & 1 & 0 \\
			0 & 1 & 0 & 1 & 1 & 0 & 1 \\
			0 & 0 & 1 & 1 & 0 & 1 & 1
		\end{array}
		\right)
	\]

	Therefore,

	\[
		X=
		\begin{pmatrix}
			1 & 1 & 1 & 0 \\
			1 & 1 & 0 & 1 \\
			1 & 0 & 1 & 1
		\end{pmatrix}.
	\]

	Now construct the parity-check matrix
	for the permuted dual code:

	\[
		H'_{C^\perp}=
		\begin{pmatrix}
			1 & 1 & 1 & 0 \\
			1 & 1 & 0 & 1 \\
			1 & 0 & 1 & 1 \\
			1 & 0 & 0 & 0 \\
			0 & 1 & 0 & 0 \\
			0 & 0 & 1 & 0 \\
			0 & 0 & 0 & 1
		\end{pmatrix}.
	\]

	Now we must undo the row permutation giving
	\[
		H_{C^\perp}=
		\begin{pmatrix}
			1 & 0 & 0 & 0 \\
			0 & 1 & 0 & 0 \\
			0 & 0 & 1 & 0 \\
			0 & 0 & 0 & 1 \\
			1 & 1 & 1 & 0 \\
			1 & 1 & 0 & 1 \\
			1 & 0 & 1 & 1
		\end{pmatrix}.
	\]

	Finally, transpose to obtain a generator
	matrix for $C$:

	\[
		G_C=H_{C^\perp}^T=
		\begin{pmatrix}
			1 & 0 & 0 & 0 & 1 & 1 & 1 \\
			0 & 1 & 0 & 0 & 1 & 1 & 0 \\
			0 & 0 & 1 & 0 & 1 & 0 & 1 \\
			0 & 0 & 0 & 1 & 0 & 1 & 1
		\end{pmatrix}.
	\]

	Now we can generate all codewords using $c=uG_C$.

	\[
		\begin{aligned}
			0000G_C & = 0000000                   \\
			1000G_C & = g_1 = 1000111             \\
			0100G_C & = g_2 = 0100110             \\
			0010G_C & = g_3 = 0010101             \\
			0001G_C & = g_4 = 0001011             \\
			1100G_C & = g_1+g_2 = 1100001         \\
			1010G_C & = g_1+g_3 = 1010010         \\
			1001G_C & = g_1+g_4 = 1001100         \\
			0110G_C & = g_2+g_3 = 0110011         \\
			0101G_C & = g_2+g_4 = 0101101         \\
			0011G_C & = g_3+g_4 = 0011110         \\
			1110G_C & = g_1+g_2+g_3 = 1110100     \\
			1101G_C & = g_1+g_2+g_4 = 1101010     \\
			1011G_C & = g_1+g_3+g_4 = 1011001     \\
			0111G_C & = g_2+g_3+g_4 = 0111000     \\
			1111G_C & = g_1+g_2+g_3+g_4 = 1111111
		\end{aligned}
	\]

	\[
		\begin{aligned}
			C=\{ &
			0000000,0001011,0010101,0011110,        \\
			     & 0100110,0101101,0110011,0111000, \\
			     & 1000111,1001100,1010010,1011001, \\
			     & 1100001,1101010,1110100,1111111
			\}.
		\end{aligned}
	\]

	Since $n=7$ and $k=4$, the number
	of cosets is

	\[
		2^{n-k}=2^3=8.
	\]

	The eight cosets can be represented by

	\[
		\begin{aligned}
			 & C+0000000  \\
			 & C+1000000  \\
			 & C+0100000  \\
			 & C+0010000  \\
			 & C+0001000  \\
			 & C+0000100  \\
			 & C+0000010  \\
			 & C+0000001.
		\end{aligned}
	\]

	Each coset contains $16$ codewords.
\end{homeworkProblem}
\begin{homeworkProblem}[2.11.2]
	\begin{enumerate}
		\item[(a)] $w=000011$

		      $000011$ belongs to the coset with leader $010000$.

		      Therefore,

		      \[
			      \begin{aligned}
				      u & = 010000        \\
				      v & = w+u           \\
				        & = 000011+010000 \\
				        & = 010011
			      \end{aligned}
		      \]
		      Meaning the decoded codeword is
		      \[
			      v=010011.
		      \]

		\item[(b)] $v=001001$

		      $001001$ belongs
		      to the coset with leader $100000$.

		      Therefore,

		      \[
			      \begin{aligned}
				      u & = 100000        \\
				      v & = w+U           \\
				        & = 001001+100000 \\
				        & = 101001
			      \end{aligned}
		      \]

		      Meaning the decoded codeword is

		      \[
			      v=101001.
		      \]

		\item[(c)] $u=001101$

		      $001101$ belongs
		      to the coset with leader $000010$.

		      Therefore,

		      \[
			      \begin{aligned}
				      u & = 000010        \\
				      v & = w+u           \\
				        & = 001101+000010 \\
				        & = 001111
			      \end{aligned}
		      \]

		      Meaning the decoded codeword is

		      \[
			      v=001111.
		      \]
	\end{enumerate}
\end{homeworkProblem}
\begin{homeworkProblem}[]


	\begin{enumerate}
		\item[(b)]
		      \[
			      C=\{0000,1010,1101,0111\}.
		      \]

		      Since $n=4$ and $k=2$, there are

		      \[
			      2^{n-k}=2^2=4
		      \]

		      possible syndromes.

		      Lets write the generator matrix $G$ as the two independent codewords

		      \[
			      G=
			      \begin{pmatrix}
				      1 & 0 & 1 & 0 \\
				      1 & 1 & 0 & 1
			      \end{pmatrix}.
		      \]
		      \[
			      R_2 \to R_2+R_1
			      \quad\Longrightarrow\quad
			      G=
			      \begin{pmatrix}
				      1 & 0 & 1 & 0 \\
				      0 & 1 & 1 & 1
			      \end{pmatrix},
			      \qquad
			      X=
			      \begin{pmatrix}
				      1 & 0 \\
				      1 & 1
			      \end{pmatrix}.
		      \]

		      so we can write the parity check matrix as

		      \[
			      H=
			      \begin{pmatrix}
				      X \\
				      I_2
			      \end{pmatrix}
			      =
			      \begin{pmatrix}
				      1 & 0 \\
				      1 & 1 \\
				      1 & 0 \\
				      0 & 1
			      \end{pmatrix}.
		      \]

		      We can start with the 0 weight error pattern

		      \[
			      0000H=00.
		      \]

		      and then we can move on to all errors of weight $1$

		      \[
			      \begin{aligned}
				      1000H & =10 \\
				      0100H & =11 \\
				      0010H & =10 \\
				      0001H & =01
			      \end{aligned}
		      \]

		      We now have all four possible syndromes totaling 4 (combine the duplicate) so we can
		      write our SDA as

		      \[
			      \begin{array}{c|c}
				      \text{Syndrome} & \text{Coset Leader} \\
				      \hline
				      00              & 0000                \\
				      01              & 0001                \\
				      10              & 1000                \\
				      11              & 0100
			      \end{array}
		      \]


		\item[(c)]

		      \[
			      C=\{00000,10100,01011,11111\}.
		      \]

		      Since $n=5$ and $k=2$, there are

		      \[
			      2^{n-k}=2^3=8
		      \]

		      possible syndromes.

		      Lets write the generator matrix $G$ again using
		      the two independent codewords:

		      \[
			      G=
			      \begin{pmatrix}
				      1 & 0 & 1 & 0 & 0 \\
				      0 & 1 & 0 & 1 & 1
			      \end{pmatrix}.
		      \]

		      Since the first two columns already form
		      the identity matrix, $G$ is in standard form:

		      \[
			      G=
			      \left(
			      \begin{array}{cc|ccc}
				      1 & 0 & 1 & 0 & 0 \\
				      0 & 1 & 0 & 1 & 1
			      \end{array}
			      \right)
			      =
			      \qquad,
			      X=
			      \begin{pmatrix}
				      1 & 0 & 0 \\
				      0 & 1 & 1
			      \end{pmatrix}.
		      \]

		      So we can write the parity-check matrix as

		      \[
			      H=
			      \begin{pmatrix}
				      X \\
				      I_3
			      \end{pmatrix}
			      =
			      \begin{pmatrix}
				      1 & 0 & 0 \\
				      0 & 1 & 1 \\
				      1 & 0 & 0 \\
				      0 & 1 & 0 \\
				      0 & 0 & 1
			      \end{pmatrix}.
		      \]

		      We can start with the weight $0$ error pattern
		      \[
			      00000H=000.
		      \]

		      And then we can move on to all errors of weight $1$:

		      \[
			      \begin{aligned}
				      10000H & =100 \\
				      01000H & =011 \\
				      00100H & =100 \\
				      00010H & =010 \\
				      00001H & =001
			      \end{aligned}
		      \]

		      Since we are still missing three syndromes,
		      we move on to error patterns of weight $2$:


		      \[
			      \begin{aligned}
				      11000H & = 111 \\
				      01100H & = 111 \\
				      00110H & = 110 \\
				      00011H & = 011 \\
				      10001H & = 101
			      \end{aligned}
		      \]

		      We now have all eight possible syndromes,
		      so we can write our SDA as

		      \[
			      \begin{array}{c|c}
				      \text{Syndrome} & \text{Coset Leader} \\
				      \hline
				      000             & 00000               \\
				      001             & 00001               \\
				      010             & 00010               \\
				      011             & 01000               \\
				      100             & 10000               \\
				      101             & 00101               \\
				      110             & 00110               \\
				      111             & 01100
			      \end{array}
		      \]
	\end{enumerate}

\end{homeworkProblem}
\end{document}
